\largerpage
\ssscp{\tsc{Irarutu-Kuri}}{11-05-02-02}
Of all the \tsc{East Bomberai} languages,
the position of Irarutu
and closely related Kuri (a.k.a.
Nabi) has remained the most unclear.
While \citet{Blust1993} linked Arguni with the other Austronesian
languages of the Bomberai peninsula,
he was unable to reach any firm position for Irarutu stating
that it “[{\ldots}] shows no known positive evidence of belonging
to the SHWNG group.
Its position for the present remains indeterminate.” \citet[31]{Kamholz2014},
assessed Irarutu as probably having “some kind of shared history
[{\ldots} with {\ldots}] SWHNG languages”
and suggested that it could be a primary branch of SHWNG.
However, at the same time he deliberately excluded it from his
investigation into the SHWNG subgroup due to its ambiguous status.
\citet[242f]{Jackson2014} presents a discussion of Irarutu historical
phonology and after considering the various hypotheses concludes that Irarutu
is probably a primary branch of CEMP which has undergone contact with SHWNG languages.
He notes that classifying Irarutu as SHWNG would probably require
positing borrowing from an unknown non-SHWNG Austronesian source.\footnote{
		\citet{SchapperZobel2024b} link Irarutu with Kowiai and
		thus within our \tsc{Eastern Islands}.
		See footnote \ref{fn:SchapperZobel2024}.}

Part of the difficulty with classifying Irarutu is due to the
heavily Papuanised nature of the language with few clear Austronesian inheritances.
In addition, words in Irarutu are usually reduced to monosyllables
by deletion of either the final or penultimate syllable thus
obscuring many Austronesian inheritances
and erasing evidence of sound changes which could be used to
group it with other languages.

Whether Irarutu and Kuri should be classified as SHWNG can be
determined by the reflexes of *j.
This is because *j
and *s consistently merge in SHWNG \citep[52]{Kamholz2014}.
We have identified six words which potentially reflect PMP *j.
Two have *j > \it{s} which is diagnostic of SHWNG,
and four have *j > \it{r} which is more consistent with \tsc{Tanimbar-Bomberai}.
Each potential reflex of PMP *j is discussed below.\footnote{
		In addition to these six words *waRəj > Irarutu \it{wara} ‘rope’
		has loss of final *j.}

\begin{enumerate}
	\item	*pajay > \it{fas} [ɸasə] ‘rice’ is probably a loan.
				Formally similar terms are widespread throughout both Austronesian
				and Papuan languages in the Bird’s Head and Bomberai Peninsula.
				As discussed in relation to \trf{11-04} on \prf{tab:11-04},
				all other Austronesian Bomberai languages (except Uruangnirin)
				have irregular retention of *p = \it{p} rather than expected
				*p > \it{f}, providing further evidence that this word has spread by contact.
				It is not surprising that the word for ‘rice’ would be a loan
				as the Bomberai region has traditionally been predominately sago oriented.
				Potential donors include Geser \it{fasa} ‘rice’ or neighbouring Wamesa \it{pas}.
				(Geser has functioned as a lingua franca in the region for many
				centuries.) \item	*bujəq > \it{fus} [ɸusə] ‘foam,
				froth’ would be unproblematic.
				But, to do due diligence,
				we note that \citet{BlustTrussel2020} reconstruct a doublet *busa
				which would regularly account for the Irarutu word.
				Furthermore, Yamdena has \it{bu{\tl}bus-e} ‘foam’.
				Thus, even if one rejects the doublet *busa as spurious (see
				footnote 15 and the discussion following \trf{11-24}),
				*j > \it{s} in the Irarutu word would be paralleled by Yamdena -- a
				\tsc{Tanimbar-Bomberai} language and could be proposed as a lexically specific innovation at the level of Proto-Tanimbar-Bomberai.
	\item	*qaləjaw > Irarutu \it{re} [re],
				\it{rre} [rəre] ‘day, sunshine’ may show *j > \it{r}, but this
				requires otherwise unattested final *-aw > \it{e}.
				\citet[262]{Jackson2014} analyses \it{rre} [rəre] as a reduplicated form \it{r{\tl}re}.
				Under this analysis, the sequence [re] would be a reflex of medial
				**lə with loss of the final syllable, including *j.
				Kuri has \it{ror} [ror] ‘day’
				and \it{rror} [rəˈrɔːr] ‘daylight’ which preserves final *j > \it{r}.
				*qaləjaw > Kuri \it{ror} requires *ə > \it{o},
				which is attested in other words.\footnote{
						Six reflexes of *ə
						have been identified in Kuri,
						of which two have *ə > \it{o}: *təlu > \it{tor} ‘three’ (Irarutu
						\it{tor}) and *dəŋəR > \it{ya-v-nogr} ‘hear’ (Irarutu \it{-f-nogr} ‘hear’).
						Other reflexes of *ə in Kuri are: *qatəluR > \it{tru} ‘egg’ (Irarutu
						\it{tru} ‘egg’), *taŋila > \it{tgra} [tə̆gra] ‘ear’ (Irarutu \it{tgra} [təŋgəra]),
						*əsa > \it{eso} ‘one’ (Irarutu \it{esu}),
						and *pusəj > \it{wuer} [wuɛrɛ] ‘navel’ (Irarutu \it{fwir} ‘navel’).}
	\item	*qapəju > Irarutu \it{-fir-} ‘bile’ is mostly unproblematic,
				though it would be only one of two potential examples of *ə >
				\it{i}, the other being *pusəj > \it{fwir} ‘navel, centre’.
	\item	*pusəj > \it{fwir} ‘navel,
				centre’ requires positing *s > {\0},
				which is not the normal reflex in Irarutu but is attested in some words.
				(Other examples of *s > {\0} include: *sida > \it{ir} ‘they,
				\tsc{3pl}’ and *nusa > \it{nu} ‘island’.) This etymology also
				requires *ə > \it{i} which would otherwise only occur in *qapəju
				> \it{-fir-} ‘bile’ discussed above.
				Kuri has ‹ʷuèrè› [wuɛrɛ] ‘navel’ (probably phonemically /wuer/
				or /uer/ with final vowel epenthesis) which requires *p > \it{w} /{\gap}\it{u}.
				This is not implausible given Kuri *p/*b > [b] {\tl} [β] {\tl}
				[f] {\tl} [v] elsewhere.
	\item	\it{wegur} [wɛŋgurə] ‘nose’ is potentially a reflex of
				PMP *ijuŋ {\tl} *ŋijuŋ {\tl} *ujuŋ.
				However, this is dependent on analysing this word as a historic compound.
				This is the analysis adopted by \citet[125]{Jackson2014} who segments it \it{-we- + -gur-}.
				Kuri has ‹ya-wegě› [yawegə̆] \it{ya-we-g} \tsc{1sg}-nose-\tsc{1sg}
				and ‹iwè› \it{i-we} \tsc{3sg}-nose \citep[34]{SmitsVoorhoeve1992a},
				thus providing comparative support for analysing the initial
				two segments of Irarutu \it{wegur} as a historically independent morpheme.
				Nonetheless, we would still need to posit irregular *i > \it{u}
				to derive **gur from the first three segments of *ŋijuŋ.
				Overall, we do not find the potential link between \it{wegur} [wɛŋgurə]
				and *(ŋ)ijuŋ convincing.
\end{enumerate}

To summarise,
if the merger of *j/*s is taken as diagnostic of SHWNG,
then placing Irarutu within this node would rest on only two
words; *bujəq > \it{fus} ‘foam’
and *pajay > \it{fas} ‘rice’ both of which are problematic: \it{fas}
‘rice’ is well attested as a loan in the region
and \it{fus} ‘foam, froth’ is potentially from the doublet *busa
‘foam’ which is also be found in Yamdena \it{bu{\tl}bus-e}.

On the other hand,
there are four potential instances of *j > \it{r} though,
admittedly, all involve problems except *qaləjaw > Kuri \it{ror} ‘day’.
Nonetheless, the explanations for the development of *qapəju
> \it{-fir-} ‘bile’ and *pusəj > Irarutu \it{fwir},
Kuri \it{(w)uer} ‘navel’ find parallels in other sound changes.
At the very worst,
there is one unproblematic instance of *j > \it{r} (*qaləjaw
> Kuri \it{ror} ‘day’)
and no unproblematic instances of *j > \it{s}.
Thus, we agree with \citet{Blust1993} that there is no positive
evidence for including Irarutu in SHWNG.

On the other hand,
there \it{is} positive evidence from the sound changes for including
Irarutu and Kuri in \tsc{Tanimbar-Bomberai}.
It shares the merger of *z/*d > **d,
as attested by *zalan ‘road’ > Fruata Irarutu \it{radni},
Kaimana Irarutu \it{rarn}, Kuri \it{rarn}
and perhaps also *zauq ‘far’ > Irarutu \it{nero},
Kuri \it{gero} (the second syllables).
In addition, Irarutu and Kuri share *j > \it{r} in four of the
six words for which *j > **r posited for \tsc{Tanimbar-Bomberai} (\trf{11-25}).

Within \tsc{Tanimbar-Bomberai}, Irarutu and Kuri share *l > \it{r},
*ŋ > \it{g}, and final vowel deletion/reduction with the \tsc{Nuclear
East Bomberai} languages (see \trf{11-30},
\trf{11-31}, and \trf{11-33}), and we therefore classify both
as members of a single \tsc{East Bomberai} node.